\cotefile{157-2}
\scanpage{243}
\mspage{189}

il faut prendre $Z$ sans élément final. \deleted{Mais}

Mais prenons $Z$ ens \unsure{ordonné}, \deleted{\unsure{et}} \unsure{avec} un objet initial $\unsure{\varphi|\emptyset|0}$ mais sans objet final, \deleted{\unsure{mais}} \inserted{et} prenons \inserted{pour} $z_0$ un élément \uline{maximal} de $Z$, de sorte que $Z'=z_0\backslash Z$ est \unsure{isom.} à $\{z_0\}$. Il faut donner \deleted{\unsure{alors}} un exemple où l'inclusion de $X_{z_0}$ dans $X$ n'est pas $\in W$. On va prendre

\diagram{\begin{tikzcd}[ampersand replacement=\&,column sep=small,row sep=tiny]
\& \deleted{\unsure{z}}a\ \deleted{\unsure{z_0}}\\
Z=\varphi \arrow[ur]\arrow[dr] \&\\
\& \deleted{\unsure{z}}b\ \deleted{\unsure{z}}=z_0
\end{tikzcd}}

donc la donnée de $X$ \unsure{cofibrée|cofibré} \unsure{scindée|scindé} sur $Z$ équivaut à un diagramme dans $\mathrm{Cat}$

\diagram{\begin{tikzcd}[ampersand replacement=\&,column sep=small,row sep=tiny]
\& X_a\\
X_0\arrow[ur,"\alpha"]\arrow[dr,"\beta"]\&\\
\& X_b\quad .
\end{tikzcd}}

Je pense que sous des conditions assez fortes (p.ex. si $\alpha,\beta$ sont des immersions ouvertes, ou des immersions \unsure{fermées} -- à vérifier !) la catégorie $X$ est isomorphe dans $\mathrm{Hot}$ à la somme amalgamée dans $\mathrm{Cat}$
\[
X_a\amalg_{X_0}X_b
\]
et il n'y \unsure{aucune} raison (sauf si on suppose $\alpha\in W$, et encore \ldots) que $X_b\unsure{\mathbin{-}|\to}X_a\amalg_{X_0}X_b$ soit dans $W$. Prenons $X_0=X_b=e$, $\beta=\mathrm{id}_e$ donc $\alpha$ est donné par la donnée d'un élément $t$ de $\mathrm{Ob}(X_a=T)$. La \deleted{catégorie $X$ est} \inserted{\unsure{somme} amalgamée} est $T$, qui n'est pas \unsure{isom.} à $e$ dans $\unsure{\mathrm{Hot}_{W}|\mathrm{Hot}_{w}}$,
% archive sheet number: 242
\scanpage{244}
\mspage{190}

si $\unsure{\overline{X}|\overline{T}}$ n'a pas été prise $W$-asphérique. Mais \deleted{on voit directement que l'\unsure{inclusion}} je vais essayer de montrer directement que l'inclusion
\[
\deleted{X_a\longrightarrow X}
\]
\deleted{est $W$-coasphérique,} $T=X_a\longrightarrow X$ est $\in W$. Pour ceci je regarde la sous-catégorie ouverte $Z''$ de $Z$, complémentaire de la sous-catégorie \unsure{pleine} fermée $Z'$, et je constate que $X''=X\mid Z''$ étant propre sur $Z''$, et $a$ étant objet final de $Z''$, \deleted{\unsure{on a}}
\[
X_a\longrightarrow X''\quad\text{dans }W,
\]
il reste à voir qu'il en est de $=$ de $X''\to X$. \deleted{Or} $X''$ est déduite de $X$ en ajoutant un unique élément \inserted{$\unsure{\xi|\zeta}$} (\deleted{\unsure{sur}} celui de $X_b$), \unsure{duquel|lequel} ne s'envoie dans aucun élément de $X''$, et qui ne reçoit aucun élément de $X_a$ -- par contre pour \deleted{\unsure{tout}} \inserted{l'unique} élément \inserted{$\xi_0$} de $X_0$ il y a une unique flèche de celui-ci dans $\xi$. On a donc une rétraction de $X$ sur $X''$, envoyant $\xi$ sur $\xi_0$ et la flèche précédente sur $\mathrm{id}_{\xi_0}$. Est-ce \unsure{un|une} \unsure{homotopisme} ? Il suffit de \unsure{prouver} \deleted{\unsure{donc}} que \unsure{endomorphisme|l'endomorphisme} \inserted{$\varphi$} de $X$, \unsure{n'est|c'est} $\sim\mathrm{id}_X$. Or on a bien $\mathrm{id}_X\leftarrow\varphi$, dont la valeur sur $X''$ est l'identité, et la valeur sur $\xi$ est l'unique $\varphi(\xi)=\xi_0\longrightarrow\xi$. C'est un peu bref, mais ça doit être vrai. (De \deleted{\unsure{toute}} façon, $X$ est la réunion de deux ouverts $X''$, et de l'ouvert $X\mid\{0,b\}=X'''$, avec $X'''\simeq\Delta^{\underline{1}}$, et $X''\cap X'''\simeq e$, donc au moins pour la coh. \inserted{à coeff. \deleted{\unsure{loc.}} \unsure{const.|crist.}} de $X$ je suis à l'aise par Mayer-Vietoris, et sais bien que c'est celle du morceau $X''$ donc celle de $X_a$, nullement celle de $e$ !
% archive sheet number: 243
\scanpage{245}
\mspage{191}

\margin{$\Vert$}
On a \deleted{donc} vu au moins ceci : \unsure{Si}\textsuperscript{*} $X$ est cofibrée donc $W$-propre sur $S$, il ne s'ensuit nullement que $\unsure{\Phi|\mathscr{D}}(X)$ soit \deleted{\unsure{cofibrée}} \inserted{même $W$-propre} sur $S^{*}$, ni que ce soit une $W$-fibrée.
% slanted
\margin{voir résultat plus précis \unsure{page suivante|pages suivantes} !}

Mais que se passe-t-il si $X\in\underline{\mathrm{Prop}}_{\unsure{\circ|0}W}S$, par exemple s'il \deleted{\unsure{provient}} vient de $\underline{\mathrm{Cofib\,sc}}_{\circ}S$ ? Alors
\deleted{\unsure{pour} \illegible{1}} \deleted{\diagram{\begin{tikzcd}[ampersand replacement=\&,column sep=small,row sep=tiny]
\&\unsure{a}\\
\unsure{0}\arrow[ur]\arrow[dr]\&\\
\&b
\end{tikzcd}}} $X$ est dans $\underline{\mathrm{Fib}}_W S$, donc
on sait déjà \unsure{que} $\Phi(X)$ est dans $\underline{\mathrm{Fib\,sc}}_{\circ W}S$, donc ce n'est pas par là qu'on fera un contre-exemple. Mais considérons un un $X$ cofibrée sur

\diagram{\begin{tikzcd}[ampersand replacement=\&,column sep=small,row sep=tiny]
\&a\\
S=0\arrow[ur]\arrow[dr]\&\\
\&b\quad,
\end{tikzcd}}

donnée par le diagramme

\diagram{\begin{tikzcd}[ampersand replacement=\&,column sep=small,row sep=tiny]
\&X_{\unsure{a|0}}\\
X_0\arrow[ur,"u"]\arrow[dr,"\unsure{v|u}"]\&\\
\&X_b
\end{tikzcd}}

dans $\mathrm{Cat}$, en supposant cette fois $u,v\in W$ ($X\in\underline{\mathrm{Cofib\,sc}}_{\circ W}S$). \deleted{\unsure{D'où $X$ est telle que}} Alors $\Phi(X)$ est la catégorie fibrée scindée sur $S$ définie par le diagramme contravariant sur $S$

\diagram{\begin{tikzcd}[ampersand replacement=\&,column sep=small,row sep=tiny]
\&X_a\arrow[dl,hook,"i"']\\
X\&\\
\&X_b\arrow[ul,hook,"j"]
\end{tikzcd}}\quad,

où $i,j$ sont les inclusions (qui sont des imm. fermées). Dire que $\unsure{\mathscr{X}|\mathfrak{X}}=\Phi(X)$ est propre sur $S$ signifie aussi que ses restrictions au dessus des \deleted{\unsure{les}} sous-catégories ouvertes $\{0,a\}$ et $\{0,b\}$ de $S$ le sont (car toute flèche $\Delta^1\to S$ se factorise par l'une ou l'autre). Cela
% archive sheet number: 244
\scanpage{246}
\mspage{192}

signifie donc aussi que l'inclusion des fibres fermées
\[
X_{\unsure{a|0}}\to X,\qquad X_b\to X
\]
sont coasphériques. Mais $X_a\to X$ ne peut être coasphérique que si $X_b=\emptyset$, car pour $x\in X_b\subset X$, $x\backslash X_a$ (ainsi que $X_a/x$) est vide, donc \uline{pas} $W$-asphérique (si on ne suppose $W=\mathrm{Fl}(\mathrm{Cat})$, cas que nous excluons). Donc $\unsure{\Phi|\mathscr{D}}(X)$ n'est propre sur $S$ \uline{que} si $X_a\deleted{\unsure{\text{et}}}\unsure{=}X_b=\emptyset$ (résultat valable pour \uline{\unsure{tte}} $X$ sur $S$, qu'elle soit cofibrée sur $S$ ou $W$-propre, ou non). Si on veut $u,v\in W$, cela implique donc que $X_0=\emptyset$, donc on a prouvé :

\margin{$\vert\!\Vert$}
Soit $X\in\mathrm{Cat}_{/S}$, $S={}$\diagram{\begin{tikzcd}[ampersand replacement=\&,column sep=tiny,row sep=tiny]
\&a\\0\arrow[ur]\arrow[dr]\&\\\&b
\end{tikzcd}}. Si $\Phi(X)$ est propre
sur $S$, alors $X=\emptyset$ (si on suppose $W\ne\mathrm{Fl}(\mathrm{Cat})$).

Mais chemin faisant est apparue une idée bizarre. \unsure{Considérons|Reconsidérons} le \deleted{\unsure{complexe}} diagramme infini, pour $X$ dans $\mathrm{Cat}_{/S}$ \deleted{$\unsure{S}$} :
\[
X\longrightarrow\Phi(X)\longrightarrow\Psi\Phi X\ \unsure{\text{---}|\to}\ \Phi\Psi\Phi(X)\longrightarrow\ldots
\]
Si $X$ est quelconque, ces flèches sont alternativement dans $W_S^d$ et $W_S^g$, donc leurs \unsure{composés|composées} sont dans $W_S$ sans plus. Mais si $X$ \deleted{\unsure{de}} \inserted{est dans} $\underline{\mathrm{Fib}}_W S$, \deleted{\unsure{alors les}} \inserted{comme} cette catégorie est \uline{stable} par $\Phi,\Psi$, et que sur $\underline{\mathrm{Fib}}_W S$ $W_S^d=W_S^g=W_S^f=W_S^{\unsure{\mathrm{univ}|\mathrm{loc}}}$, on voit que dans ce cas les flèches sont toutes dans $W_S^f$. \deleted{En tous cas on} On a envie maintenant de poser $\Phi_\infty(X)=\underrightarrow{\lim}$ de tout ça
% archive sheet number: 245
\scanpage{247}
\mspage{193}

et en définissant symétriquement
\[
\Psi_\infty(X)=\underrightarrow{\lim}\bigl(X\to\Psi(X)\longrightarrow\unsure{\Phi|\mathscr{D}}\Psi(X)\longrightarrow\ldots\bigr).
\]
Pour que ce soit utilisable, il faut poser l'axiome

\uline{$W(6\ \mathrm{bis})$} (comparer $\overline{\underline{\mathrm{XII}}}$, pages $14,15$) La limite inductive d'une suite de morphismes $f_i:X_i\unsure{\mathbin{-}|\to}Y_i$ \unsure{i.e.|s.} pour un morphisme de systèmes inductifs dans $\mathrm{Cat}$ indexés par $\unsure{\mathbb N|N}$, est dans $W$ si les $f_i$ le sont.

Cet axiome est voisin de $W(6)$ de loc. cit., mais n'est impliqué par lui que moyennant $W(4\mathrm b)$). Il est plus fort dans le cas de syst. inductifs indexés par $N$ (ou aussi, plus gén., de syst. inductifs sur des ens. ordonnés filtrants ayant un ens. cofinal dénombrable), mais ne dit rien sur les syst. inductifs sur des ens. ordonnés d'indices plus généraux.

Cet axiome implique aussi qu'une limite inductive filtrante \deleted{dans $\underline{\mathrm{Cat}}_{/S}$} \inserted{de \unsure{morphismes|flèches}} \deleted{\unsure{est}} $f_i$ est $W$-lisse, resp. $W$-propre, si les $f_i$ le sont. \deleted{Ainsi, \unsure{pour}} On trouve alors que $\Phi_\infty(X)$ est à la fois $W$-lisse (comme limite inductive des $\Phi(X),\Phi\Psi\Phi(X)\ldots,$ $\Phi(\Psi\Phi)^n(X)$ \unsure{etc.}, \deleted{\unsure{cha}} \unsure{lesquelles|lesquels|quelles} sont $W$-lisses sur $S$) et $W$-propre (comme limite inductive des $\Psi\Phi(X),\ldots(\Psi\Phi)^n(X)\ldots$). \deleted{\unsure{Donc}} D'autre part, $X\longrightarrow\Phi_\infty(X)$ est dans $W_S^f$ -- car \inserted{\unsure{$W_S^f$ est stable par}} \deleted{\unsure{une}} limites inductives \inserted{filtr.} dénombrables dans $\mathrm{Cat}_{/S}$. On a donc trouvé un foncteur
% archive sheet number: 246
\scanpage{248}
\mspage{194}

\[
\unsure{\Phi|\mathscr{D}}_\infty:\underline{\mathrm{Fib}}_W S\xrightarrow{\deleted{\unsure{\Phi}}}\underline{\mathrm{Perf}}_W S\subset\underline{\mathrm{Fib}}_W S
\]
et un morphisme fonctoriel
\[
i_X:X\xrightarrow{\deleted{\unsure{i_X}}}\Phi_\infty(X),\qquad
\tikz[baseline=(M.base)]\node[draw,ellipse,inner sep=3pt](M){$i_X\in W_S^f$};\ \text{si }X\in\underline{\mathrm{Fib}}_W S
\]

Ce foncteur permet alors, en l'absence de $W(4\,\mathrm b)$, de faire ce que le foncteur $C$ refusait \unsure{obstinément|absolument}. (Il est vrai que $C$ a l'avantage d'être défini sur $\mathrm{Cat}_{/S}$ tout entier, et d'y être \unsure{vraisemblablement|raisonnablement} $=$ fort utile, quand on a $W(4\mathrm b)$. Mais $\Phi_\infty$ est, lui aussi, défini sur $\mathrm{Cat}_{/S}$ \deleted{tout} entier, ainsi que $i_X$
\[
\begin{aligned}
\Phi_\infty:\mathrm{Cat}_{/S}&\longrightarrow\underline{\mathrm{Perf}}_W S\\
i_{\unsure{S|X}}:\mathrm{id}&\to\Phi_\infty\qquad\text{i.e. }
\begin{cases}
i_X:X\longrightarrow\Phi_\infty(X),\\
i_X\in W_S
\end{cases}
\end{aligned}
\]
avec le seul handicap (inévitable) que si \deleted{$X$ \unsure{n'est}} $X$ n'est pas dans $\underline{\mathrm{Fib}}_W S$, on ne peut dire mieux que $i_X\in W_S$, à l'exclusion de $i_X\in W_S^f$. \deleted{Pourtant, il suffit de} Mais pour $C(X)$ c'était \unsure{forcément|précis.} pareil. Donc on a l'impression que $\Phi_\infty$ (ou au choix, $\Psi_\infty$) bonnet blanc, blanc bonnet, joue, \deleted{\unsure{le}} en présence de $W(6\ \mathrm{bis})$, le même rôle que $C$ en présence de $W(4\mathrm b)$.
% slanted
\margin{\uline{En fait} $\Phi_\infty=C$, si $C$ est défini \unsure{par} des chemins infinis d'\unsure{1|un} seul côté \unsure{seulement} ; si \unsure{$C$ défini|je définis} $C$ \unsure{via|avec} les chemins \uline{finis} de \unsure{long.} \unsure{quelconque} \unsure{des} deux \unsure{côtés}, il prend en effet ses \unsure{val.} \unsure{ds} $\underline{\mathrm{Perf}}_W S$.}
De plus (en présence des deux axiomes à la fois) le foncteur $C$ a-t-il l'avantage esthétique de prendre ses valeurs dans $\underline{\mathrm{Perf}}_W S$, au lieu de $\underline{\unsure{\mathrm{P}|\mathrm{Pf}}}_W S$ seulement. Mais il s'impose de
% archive sheet number: 247
\scanpage{249}
\mspage{\unsure{195|175}}

se rendre compte à présent que $\unsure{\Phi|\mathscr{D}}_\infty$ a un avantage exactement du $=$ ordre sur $C$, en introduisant
\[
\begin{aligned}
W_{\unsure{\omega|w}}''={}&\text{plus petite partie }\Sigma\text{ de }\mathrm{Fl}(\mathrm{Cat})\\
&\text{satisfaisant }W(1,2\ \mathrm{bis},3,6\ \mathrm{bis})
\end{aligned}
\]
et noter que
\[
\Phi_\infty(X)\in\mathrm{Ob}\ \underline{\mathrm{Perf}}_{W_\omega''}S,\qquad
\left|\begin{aligned}
i_X&\in\bigl(W_\omega''\deleted{\unsure{\bigr)_S^f}}\bigr.\\[-2pt]
&\deleted{\unsure{(W_\omega'')_S^f}}
\end{aligned}\right.
\]
-- mais là je m'aperçois qu'on a
\[
W_\omega^{\unsure{\prime\prime\prime|\prime|b}}\subset W_\omega''
\]
donc $C$ est quand même un peu meilleur à cet égard \ldots

Quand on dispose de $\Phi_\infty$ avec les propriétés qu'on vient d'expliciter, on trouve que \deleted{\unsure{les}} l'inclusion
\[
\underline{\mathrm{Perf}}_W S\to\mathrm{Fib}_W(S)
\]
% slanted
\margin{On peut y remplacer $=$ $\underline{\mathrm{Perf}}_W S$ par $\underline{\mathrm{Perf}}_{W_\omega''}S$, \uline{ou} par toute \inserted{autre} sous-cat. pleine de $\underline{\mathrm{Fib}}_W(S)$ contenant $\underline{\mathrm{Perf}}_{W_\omega''}S$.}
induit une équivalence non seulement pour la localisation par $W_S$, mais aussi pour les localisés fins
\[
\begin{aligned}
(\underline{\mathrm{Perf}}_W S)(W_S^f)^{-1}&\longrightarrow(\mathrm{Fib}_W(S))(W_S^f)^{-1}.\\
&\qquad\overset{\mathrm{def}}{=}\mathrm{HOT}_W^{\unsure{\mathrm{lc}|\mathrm{le}|\mathrm{loc}}}(S)
\end{aligned}
\]

Mais on ne peut \deleted{\unsure{pour}} autant réduire le nombre des \unsure{localisés|localisions} \unsure{utiles} à trois, en l'absence de $W(4\mathrm b)$. Tout au plus, ayant un foncteur $\Phi_\infty$ qui est plus souple que $C$, en l'absence de $W(4\mathrm b)$, du fait qu'on sait de toute façon que
\[
i_X:X\longrightarrow\Phi_\infty(X)
\]
est dans $W_S^f$ si $X$ est dans $\underline{\mathrm{Fib}}_W S$, peut-on
% archive sheet number: 248
\scanpage{250}
\mspage{196}

espérer s'en tirer avec \deleted{la} variante affaiblie de $W(4\mathrm b)$, obtenue en se bornant dans l'énoncé de $W(4\mathrm b)$ (p. 171) au cas où $X,Y$ sont dans $\underline{\mathrm{Perf}}_{W_{\unsure{\omega|w}}''}S$. Cette fois, ça suffit en effet à impliquer que le foncteur
\[
\lambda:(\underline{\mathrm{Fib}}_W S)(W_S^f)^{-1}\to(\underline{\mathrm{Fib}}_W S)(W_S)^{-1}
\]
est une équivalence de catégories, puisque on peut le ``calculer'' aussi bien en utilisant $\underline{\mathrm{Perf}}_{W_\omega''}S$ au lieu de $\underline{\mathrm{Fib}}_W$, or \deleted{\unsure{cette fois}} l'hyp. $W(4\ \mathrm b(W_\omega''))$ implique que sur $\underline{\mathrm{Perf}}_{W_\omega''}S$ les deux localiseurs \inserted{$W_S^f,W_S$} sont \uline{égaux}.

Il faut une conclusion pratique utilisable de toutes ces perplexités et observations. Essayons

\uline{Proposition} Supposons que l'on soit sous l'une des conditions suivantes

$1^\circ)$ L'axiome $W(4\mathrm b)$ est satisfait dans le cas de catégories $X,Y$ \deleted{\unsure{qui se}} qui sont dans $\underline{\mathrm{Fib\,sc}}_{\circ W}S$. (Cette condition \inserted{sur $W$} est en fait \unsure{autoduale|autodual.}.)
% slanted
\margin{mais c'est \uline{équivalent} à $W(4\mathrm b)$ ! cf $\overline{\underline{\mathrm{XVI}}}$ 4.4. loc. cit.}

$2^\circ)$ L'axiome $W(4\mathrm b)$ est satisfait dans le cas de catégories $X,Y$ qui sont dans $\underline{\mathrm{Perf}}_{W_\omega''}S$, où $W_\omega''$ est la plus petite ens. $\Sigma\subset\mathrm{Fl}(\mathrm{Cat})$ satisfaisant les conditions $W(1,2\ \mathrm{bis},3,6\ \mathrm{bis})$. (On peut aussi se contenter de l'a\deleted{\unsure{x}} condition, plus forte mais plus \unsure{compréhensible}, que $W(4\mathrm b)$ est OK si $X,Y$ sont dans $\underline{\mathrm{Perf}}_W S$ !) De plus, on suppose que $W$ satisfait $W(6\ \mathrm{bis})$.
% slanted
\margin{mais cela implique aussi $W(4\mathrm b)$ cf loc. cit.}

Sous ces conditions, le foncteur \inserted{\unsure{can.}}
\[
\mathrm{HOT}_W^{\unsure{\mathrm{lc}|\mathrm{le}}}(S)\to\mathrm{HOT}_W^{\unsure{\mathrm{lc}!|\mathrm{le}!}}(S)
\]
est une équiv. de cat., (cf diagramme p. \unsure{178|170}), donc $(\unsure{\rho|g},i);(\unsure{\sigma|\rho},j)$ \deleted{sont} \inserted{\unsure{peuvent} être regardés comme} des couples de foncteurs entre deux catégories (cf \unsure{(15)} p. 181). \uline{Ce sont des couples de foncteurs adjoints}, donc $i,j$ sont pl. fidèles, et on a les formules d'adjonction \unsure{(16)}, page 182.
% archive sheet number: 249
\scanpage{251}
\mspage{197}

Pour \deleted{la} prouver \deleted{\unsure{le fait que soit un foncteur}} \inserted{l'adjonction des $(\unsure{\rho|g},i)$ (ou dualement des $(\sigma,j)$)}, \deleted{\unsure{on utilise}} on \deleted{\unsure{doit}} exprime \inserted{les} deux catégories en présence $\mathrm{HOT}_W(S),\mathrm{HOT}_W^{\unsure{\mathrm{lc}|\mathrm{le}}}(S)$ comme $(\mathrm{Cat}_{/S})(W_S^d)^{-1}$, $(\mathrm{Cat}_{/S})\deleted{\unsure{^{-1}}}(W_S)^{-1}$, le foncteur $\rho$ étant la localisation canonique allant du premier vers le second, et le foncteur $i$ étant donné par un foncteur d'\unsure{inversion|immersion}
\[
\unsure{\gamma|r}:(\mathrm{Cat}_{/S},W_S)\to(\mathrm{Cat}_{/S},\overline{W_S^{\unsure{d|1}}})
\]
qui est $C$ dans le cas $1^\circ$, $\unsure{\Phi|\mathscr{D}}_\infty$ dans le cas $2^\circ$.

Les morphismes d'adjonction
\[
\rho i\longrightarrow\mathrm{id},\qquad i\rho\longleftarrow\mathrm{id}
\]
doivent se décrire alors comme dans le cas où on suppose \deleted{tout} bonnement $W(4\mathrm b)$, et les compatibilités se voir de même : on est ramené à prouver que deux flèches

\diagram{\begin{tikzcd}[ampersand replacement=\&]
\gamma(X)\arrow[r,shift left,"{i_{\gamma(X)}}"]\arrow[r,shift right,"{\gamma(i_X)}"']\&\gamma(\gamma(X))
\end{tikzcd}}

sont égales dans $(\mathrm{Cat}_{/S})(W_S^d)^{-1}$. La vérification dans le cas $1^\circ)$, où $\gamma=C$, ne doit pas être \unsure{changée}. Dans le cas $2^\circ)$ où $\gamma=\Phi_\infty$, je ne l'ai pas faite. Je ne la ferai que sous la pression du besoin.

\textcircled{6} \uline{Loi contravariante} \uline{de} $\mathrm{HOT}_W(S),\mathrm{HOT}_W^{\mathrm{lc}}(S)$

Il y a \deleted{pour $\unsure{f}$}
\[
f:S'\unsure{\mathbin{-}|\to}S
\]
un hom. \unsure{des|de} diagrammes

\diagram{\begin{tikzcd}[ampersand replacement=\&,row sep=small]
\&\mathrm{HOT}_W^{\mathrm{lc}}(S)\arrow[dl]\arrow[dr]\&\\
\mathrm{HOT}_W(S)\&\&\mathrm{HOT}_W^{\unsure{\circ|0}}(S)
\end{tikzcd}}
% archive sheet number: 250
\scanpage{252}
\mspage{198}

dans le diagramme similaire pour $S$, par des foncteurs $f^*$ qui sont définis respectivement sur
\[
\begin{aligned}
\underline{\mathrm{Liss}}_W S&\quad\text{pour }\deleted{\unsure{\mathrm{H}}}\mathrm{HOT}_W(S)\to\mathrm{HOT}_W(S')\\
\underline{\mathrm{Prop}}_W S&\quad\text{pour }\mathrm{HOT}_W^{\unsure{\circ|0}}(S)\ \unsure{\text{---}|\to}\ \mathrm{HOT}_W^\circ(S')\\
\underline{\mathrm{Fib}}_W S&\quad\text{pour }\mathrm{HOT}_W^{\unsure{\mathrm{lc}|\mathrm{le}}}(S)\to\mathrm{HOT}_W^{\mathrm{lc}}(S'),
\end{aligned}
\]
par l'intermédiaire des foncteurs
\[
f^*:X\ \unsure{\text{---}|\to}\ X\times_S S'\quad
\left|\begin{aligned}
\underline{\mathrm{Liss}}_W S&\to\underline{\mathrm{Liss}}_W S'\\
\underline{\mathrm{Prop}}_W S&\to\underline{\mathrm{Prop}}_W S'\\
\underline{\mathrm{Fib}}_W S&\to\mathrm{Fib}_W S'
\end{aligned}\right.\quad,
\]
lequel foncteur est compatible avec les localiseurs pertinents, qui sont $W_S^f$ dans tous les cas. Par contre, en l'absence d'hypothèses draconiennes sur $W$, \inserted{ou sur $S$,} on ne peut dire que $f^{\unsure{*|d}}$ soit compatible avec $W_S,W_{S'}$ sur une quelconque des catégories du diagramme p. 175, permettant de définir \deleted{\unsure{les}} $\mathrm{HOT}_W^{\mathrm{lc}!}(S)$, de façon à nous donner un foncteur
\[
\mathrm{HOT}_W^{\mathrm{lc}!}(S)\longrightarrow\mathrm{HOT}_W^{\mathrm{lc}!}(S').
\]
En l'absence d'axiome sur $W$, le seul cas qui me soit connu est celui où $f$ est \uline{une $W$-équivalence universelle}, i.e. une $W$-fibration à fibres $W$-asphériques, auquel cas il est
% slanted
\margin{Ça s'appelle une \uline{$W$-fibration} \uline{triviale}.}
évident que pour \deleted{\uline{toute}} flèche \deleted{\unsure{$f$ de}} $u:X\to Y$ dans $\mathrm{Cat}_{/S}$, $u\in W_S$ \deleted{\unsure{ssi}} $(u':X'\to Y')\in W_{S'}$. On trouve alors un hom. du diagramme complet (\unsure{de|de p.} \deleted{\unsure{p.}} \unsure{176|175} pour $S$, dans le diagramme correspondant pour $S'$.

Si on suppose $W(4\mathrm b)$, alors $W_S=W_{\unsure{S|g}}^f$ sur $\underline{\mathrm{Fib}}_W S$, donc sur cette catégorie-là le foncteur $f^*$ est compatible (sans mérite !) avec les localiseurs $W_S,W_{S'}$, mais ce n'est pas nouveau vraiment
% archive sheet number: 251
\scanpage{253}
\mspage{199}

p.r. à ce qui a été dit plus haut pour $\underline{\mathrm{Fib}}_W S,\underline{\mathrm{Fib}}_W S'$ et pour les $\mathrm{HOT}_W^{\unsure{\mathrm{lc}|\mathrm{le}}}$.

Mais si on suppose \uline{$W(4\,\unsure{\mbox{-}|\mathrm a|\cdot})$}, et \uline{$f\in\mathrm{Fib}_W$}, alors $f^*$ sur $\mathrm{Cat}_{/S}$ tout entier est compatible avec les localiseurs $W_S,W_{S'}$, et on trouve un calcul \deleted{\unsure{fonctoriel}} de $f^*$ sur $\mathrm{HOT}_W^{\mathrm{lc}!}(S)\approx\mathrm{HOT}_W^{\mathrm{lc}}(S)$ par
\[
f^*\bigl(\mathrm{hot}^{\mathrm{lc}}_{S,W}(X)\bigr)=\mathrm{hot}^{\mathrm{lc}}_{S',W}(X\times_S S').
\]
Donc c'est là un cas, exceptionnellement, où les foncteurs $f^*$ pour les $\mathrm{HOT}_W$ sont compatibles non seulement \deleted{\unsure{pour les}} \inserted{aux foncteurs d'}inclusion $i,j$ (de la page 176), mais aussi aux foncteurs de localisation $\unsure{\rho|g},\sigma$.

\margin{\uline{Oubli} : Si $p:S'\to S$ est $W$-propre, alors $p^*:\mathrm{Cat}_{/S}\ \unsure{\text{---}|\to}\ \mathrm{Cat}_{/S'}$ est compatible avec $W_S^d,W_{S'}^d$, et $p^*:\mathrm{HOT}_W(S)\to\mathrm{HOT}_W(S')$ peut se ``calculer'' par la formule $p^*(\mathrm{hot}_{S,W}(X))=\mathrm{hot}_{S',W}(X\times_S S')$ pour \uline{tt} $X$ dans $\mathrm{Cat}_{/S}$.}

\textcircled{7} \uline{Produit \deleted{\unsure{de}} dans} $\mathrm{HOT}_W(S)$

\deleted{\unsure{Dans}} le foncteur
\[
(X,Y)\longmapsto X\times_S Y:\ \underline{\mathrm{Liss}}_W S\times\underline{\mathrm{Liss}}_W S\ \unsure{\text{---}|\to}\ \underline{\mathrm{Liss}}_W S
\]
est compatible avec les localiseurs $W_S^f\times W_S^f$ et $W_S^f$, \unsure{donc|donne} \unsure{passe|passage} aux \unsure{localisées|localisés} \unsure{un|en}
\[
\left\{\begin{aligned}
\mathrm{HOT}_W(S)\times\mathrm{HOT}_W(S)&\to\mathrm{HOT}_W(S)\\
(\unsure{\xi|\zeta},\unsure{\eta|y})&\longmapsto\xi\times_S\eta .
\end{aligned}\right.
\]
% slanted
\margin{\uline{NB} Si $X,Y$ sont \uline{\unsure{quelc.}} dans $\mathrm{Cat}_{/S}$, on a $\mathrm{hot}_{S,W}(X)\times_S\mathrm{hot}_{S,W}(Y)=\mathrm{hot}_{S,W}(\unsure{\Phi|\mathscr{D}}_S(X)\times_S\Phi_S(Y))$. Si $X$, ou $Y$ lisse \deleted{\unsure{sur}} $S$, on peut y remplacer $\Phi X$, resp $\Phi(Y)$ par $X$ resp. $Y$.}
Ce produit est associatif, unitaire, comm. avec \deleted{toutes} les compatibilités habituelles, comme il résulte du fait que c'est \unsure{comme ça|connu si.} \unsure{dans|ds} $\underline{\mathrm{Liss}}_S$. L'élément unité est la classe de $S\xrightarrow{\mathrm{id}}S$. La sous-catégorie $\mathrm{HOT}_W^{\mathrm{lc}}(S)$ est stable par le produit, qui induit
% archive sheet number: 252
\scanpage{254}
\mspage{200}

\[
\mathrm{HOT}^{\unsure{\mathrm{lc}|\mathrm{le}}}_W(S)\times\mathrm{HOT}^{\mathrm{lc}}_W(S)\to\mathrm{HOT}^{\mathrm{lc}}_W(S)
\]
lui-même induit par le foncteur \inserted{produit} $(X,Y)\mapsto X\times_S Y$ dans $\underline{\mathrm{Fib}}_W S$.

Si $X$ est quelconque dans $\mathrm{Cat}_{/S}$, le foncteur
\[
\mathrm{mult}_{X/S}:Y\to X\times_S Y:\ \underline{\mathrm{Liss}}_W S\to\mathrm{Cat}_{/S}
\]
est compatible avec les localiseurs $W_S^d$, car ce foncteur peut s'écrire comme le composé du foncteur correspondant à \uline{val. ds} \deleted{$\mathrm{Cat}_{/X}$} $\underline{\mathrm{Liss}}_W X$, lequel est compatible avec les localiseurs $W_S^d,W_X^d$, et le foncteur ``oubli de $X$'' de $\underline{\mathrm{Liss}}_W X$ (\deleted{ou} $=$ de $\mathrm{Cat}_{/X}$) dans $\mathrm{Cat}_{/S}$, lequel est compatible avec $W_X^d$ et $W_S^d$. On trouve ainsi un foncteur
\[
\left\{\begin{aligned}
\mathrm{mult}_{X/S}:&\quad\mathrm{Hot}_W(S)\to\mathrm{Hot}_W(S)\\
\mathrm{mult}_{X/S}(\mathrm{hot}_{S,W}\deleted{\unsure{\mathrm{hot}}}(Y))&\overset{\mathrm{def}}{=}\mathrm{hot}_{S,W}(X\times_S Y),\ \text{si }Y\ \underline{\text{lisse}}\text{ sur }S^{*}
\end{aligned}\right.
\]
% slanted
\margin{Cette formule est également valable si $Y$ est dans $\underline{\mathrm{Fib}}_W S$, cf $(*)$ p. 202.}
\unsure{puisqu'ici|puisque} \deleted{\unsure{ici}} $\mathrm{HOT}_W(S)$ peut s'obtenir comme localisé tant de $\underline{\mathrm{Liss}}_W S$ que de $\mathrm{Cat}_{/S}$, par $W_S^d$. Quand $X$ est lui-$=$ \inserted{$W$-}lisse sur \deleted{\unsure{sa}} $S$, ce foncteur n'est autre que
\[
\left\{\begin{aligned}
\mathrm{mult}_{\unsure{\xi|\zeta}}:&\quad\unsure{\eta|y}\longmapsto\xi\times_S\eta,\\
&\text{où }\xi=\mathrm{hot}_{S,W}(X).
\end{aligned}\right.
\]
Mais il n'en est pas de $=$ si $X$ n'est pas lisse, i.e. il faut distinguer dans ce cas, même \inserted{pour $Y$ \unsure{$W$-|lui-$=$ }lisse sur $S$}
\[
\mathrm{hot}_{S,W}(X\times_S Y)\quad\text{et}\quad\mathrm{hot}_{S,W}(X)\times_S\mathrm{hot}_{S,W}(Y).
\]
\deleted{\unsure{Pour $X$ et $Y$ qui}} \inserted{Le deuxième terme peut} s'expliciter comme $\mathrm{hot}_{S,W}(\unsure{\Phi|\mathscr{D}}(X)\times_S Y))$ qui reçoit $\mathrm{hot}_{S,W}(X\times_S Y)$, i.e.
% archive sheet number: 253
\scanpage{255}
\mspage{201}

on a un hom can.
\[
\begin{aligned}
\underline{\mathrm{hot}_{S,W}(X\times_S Y)}&\longrightarrow
\underline{\mathrm{hot}_{S,W}(X)\times_S\mathrm{hot}_{S,W}(Y)},\\
\Vert&\hspace{7em}\Vert\\
\mathrm{mult}_{X/S}(Y)&\qquad\mathrm{hot}_{S,W}(\unsure{\Phi|\mathscr{D}}(X)\times_S Y)
\end{aligned}
\]
induit par
\[
X\times_S Y\longrightarrow\Phi(X)\times_S Y
\]
mais cette\deleted{\unsure{ flèche}} dernière flèche n'a aucune raison d'être dans $W_S^d$, \deleted{\unsure{donc}} et n'induit \inserted{donc} pas \unsure{néc.$^{t}$} un iso dans $\mathrm{HOT}_W(S)$. (Contrexemple page 154.)

Il faut regarder à part le cas \uline{où $X$ est dans $\underline{\mathrm{Fib}}_W S$}, j'ai envie de dire que dans ce cas on a
\[
\begin{aligned}
\mathrm{hot}_{S,W}(X\times_S Y)&\overset{\sim}{\longrightarrow}
\underline{\mathrm{hot}_{S,W}(X)\times_S\mathrm{hot}_{S,W}(Y)}^{*}\\
&\hspace{5em}\Vert\\
&\quad\mathrm{hot}_{S,W}(\Phi(X)\times_S\Phi(Y))
\end{aligned}
\]
% slanted
\margin{Cette formule n'est établie qu'en supposant ou bien $Y$ $W$-lisse (en plus de $X$ $W$-fibré sur $S$ !), ou bien $X$ $W$-propre, p. ex. $X$ $W$-parfait \unsure{?|!}.}
pour \uline{tout} $Y$ dans $\mathrm{Cat}_{/S}$ \inserted{(pas néc. $W$-lisse)}, a fortiori que dans ce cas
\[
\left\{\begin{aligned}
\mathrm{mult}_{X/S}&=\mathrm{mult}_{\unsure{\xi|\zeta}}\qquad\text{dans }\mathrm{HOT}_W(S)\\
\text{où }\xi&=\mathrm{hot}_{S,W}(X/S),
\end{aligned}\right.\quad\bigl(\text{si }X\text{ dans }\underline{\mathrm{Fib}}_W\unsure{S|S'}\bigr)\ \Vert\tag*{$(**)$}
\]
\deleted{\unsure{mais il faut prouver que} \illegible{1}} \deleted{Quand} \inserted{Dans le cas où} $Y$ est lisse sur $S$ (ce qui suffit pour établir $(**)$), il suffit de \deleted{\unsure{voir}} \inserted{se rappeler} que $X\to\Phi(X)$, \deleted{étant dans} \deleted{$W_S^{\unsure{f}}$} dans $W_S^f$ (résulte du fait que $X$ est $W$-fibré sur $S^{*}$, donc $X_s\unsure{\mathbin{-}|\to}s\backslash X$ est dans $W$, $\forall s\in S$),
% slanted
\margin{il suffit $=$ $W$-prélisse sur $S$}
donc $\underbrace{X\times_S Y}_{X_Y}\to\underbrace{\Phi(X)}_{\Phi(X)_Y}\times_S Y$ est aussi dans $W_Y^f$, \deleted{et comme les deux sont $W$-prélisses sur $S$,} donc (comme \deleted{\unsure{$X,Y\in$}} $X_Y,\Phi(X)_Y$ sont $W$-fibrés sur $Y$) dans $W_Y^d$, donc aussi dans $W_S^d$. En fait, $=$ sans supposer $Y$ lisse, ceci montre que l'on a
% archive sheet number: 254
\scanpage{256}
\mspage{202}

\[
\mathrm{hot}_{S,W}(X\times_S Y)\overset{\sim}{\longrightarrow}\mathrm{mult}_{Y/S}(\mathrm{hot}_{S,W}(X))\quad\left|\ \begin{aligned}&\text{si }X\text{ est}\\&\text{dans }\underline{\mathrm{Fib}}_W S.\end{aligned}\right.\tag*{$(*)$}
\]
Pour avoir ceci pour $Y$ quelconque, on peut par ce qui précède supposer $X$ lisse (quitte à le remplacer par $\unsure{\Phi|\mathscr{D}}(X)$), \unsure{donc $\alpha$ un|donc $\in$ à} $\underline{\mathrm{Liss}}_{\unsure{\circ|0}W}S$. Il faut voir que
\[
X\times_S Y\unsure{\mathbin{-}|\longrightarrow}X\times_S\Phi(Y)
\]
est \deleted{\unsure{dans}} \unsure{une|un} dans $W_S^d$, \deleted{\unsure{supp}} Or \deleted{\unsure{que}} si $X$ est \uline{$W$-propre} sur $S$, le chgt. de base $X\unsure{\mathbin{-}|\to}S$ est compatible avec $W_S^d,W_X^d$ (sur $\mathrm{Cat}_{/S}$ \unsure{tout} entier), et on gagne. (Ça doit être faux sans hypo. de $W$-propreté sur $X$, \uline{ou} de $W$-\unsure{lissité} sur $Y\ldots$)

Quand on a
\[
f:S'\unsure{\mathbin{-}|\longrightarrow}S
\]
alors
\[
f^{\unsure{*|1}}:\mathrm{HOT}_W(S)\unsure{\mathbin{-}|\longrightarrow}\mathrm{HOT}_W(S'),
\]
est compatible avec les structures de \unsure{produit|produits}~!

\textcircled{8} Les foncteurs $f_!$ : \unsure{Théorèmes|Théorème} d'adjonction, propriétés du dérivateur $\mathrm{HOT}_W$.

Soit
\[
f:S'\unsure{\mathbin{-}|\longrightarrow}S
\]
dans $\mathrm{Cat}$, alors le foncteur induit
\[
f_!:\mathrm{Cat}_{/S'}\longrightarrow\mathrm{Cat}_{/S}
\]
est compatible avec les localiseurs $(W_{S'}^d,W_S^{\unsure{1|d}})$, ou $(W_{S'}^g,W_S^g)$, \unsure{ou} $(W_{S'},W_S)$. On trouve donc un foncteur induit pour les localisées, et \unsure{diagramme} \unsure{comm.}
% archive sheet number: 255
\scanpage{257}
\mspage{203}

\diagram{\begin{tikzcd}[ampersand replacement=\&]
\mathrm{HOT}_W(S')\arrow[dd,"f_!"']\arrow[dr,"\rho'"] \& \& \mathrm{HOT}^{\unsure{\circ|0}}_W(S')\arrow[dd,"f_!^{\circ}"]\arrow[dl,"\sigma'"']\\
\& \mathrm{HOT}^{\unsure{\mathrm{lc}|\mathrm{le}}!}_W(S')\arrow[dd,"f_!^{\mathrm{lc}!}"]\&\\
\mathrm{HOT}_W(S)\arrow[dr,"\rho"]\&\&\mathrm{HOT}^{\circ}_W(S)\arrow[dl,"\sigma"']\\
\&\mathrm{HOT}^{\mathrm{lc}!}_W(S)\&
\end{tikzcd}}

On n'a pas de foncteur \inserted{$f_!^{\mathrm{lc}}$} au niveau des $\mathrm{HOT}^{\mathrm{lc}}_W$, sauf si on fait une hypothèse, telle \uline{$W(4b)$}, qui assure que $\mathrm{HOT}^{\mathrm{lc}!}_W\simeq\mathrm{HOT}^{\mathrm{lc}}_W$.

Le théorème principal, que \deleted{les deux} \unsure{foncteur|foncteurs} \deleted{\unsure{sont adj. à gauche de}} \inserted{$f_!$ est adjoint à gauche} de $f^*$, est valable \uline{sous la seule hypothèse} $W(1,2,3)$.\textsuperscript{*}
% slanted
\margin{\unsure{i.e.} avec les nouvelles notations \deleted{$\mathrm{loc}(1,\unsure{2\,\mathrm{bis}},3)$} $\mathrm{loc}(1,\unsure{2|?},3\,\mathrm{bis})$ cf $\overline{\underline{\mathrm{XVI}}}$)}
Donc de même pour $f_!^\circ,f_\circ^*$. Mais comme on n'a pas de foncteur \deleted{\unsure{sur}} $f^*_{\mathrm{lc}!}$, la question d'adjonction ne se pose pas pour $f$ général, et sans hyp. $W(4b)$ sur $W$. Elle ne se pose, apparemment, que pour $f$ un $W$-fibré trivial (cf fin de \textcircled{6}, p.~\unsure{198|195}), et dans ce cas l'adjonction ne doit pas faire de difficulté. Mais tant que je n'ai pas de sentiment bien net au sujet de la signification des $\mathrm{Hot}^{\mathrm{lc}!}_W$, ce n'est pas là un résultat \unsure{bien|très} \unsure{bandant|transcendant|intéressant}.

L'adjonction de $f_!,f^*$ pour $\mathrm{HOT}_W$ permet déjà de \unsure{montrer} que le prédérivateur $\mathrm{HOT}_W$
% archive sheet number: 256
\scanpage{258}
\mspage{204}
% archive sheet number: 257

satisfait les conditions $\mathrm{Der}\,1,\mathrm{Der}\,4',\mathrm{Der}\,5'$ des dérivateurs. \deleted{\unsure{Der 2}} \unsure{plus|De plus}, on a \unsure{ce} qu'on peut appeler la propriété $\mathrm{Der}\,3b$ pour $\mathrm{HOT}_W$ en tant que $W$-dérivateur, i.e. \uline{Si} $f\in W$, \uline{alors} $f$ \uline{est une} $\mathrm{HOT}_W$-\uline{équivalence}\textsuperscript{*}
% slanted
\margin{On le prouve en interprétant la notion de $\mathrm{HOT}_W$-équiv. par la $\mathrm{HOT}_W$-\uline{homotopie}. Cela revient finalement à ceci : si $f:S'\unsure{\mathbin{-}|\to}S$ est \unsure{ds|dans} $W$, alors $\forall T$, $S'\times T\unsure{\mathbin{-}|\to}S\times T$ \unsure{aussi}~!}
(\unsure{moyennant} les \deleted{seuls} axiomes $W(1,2,3)$ sur $W$). \deleted{Plus} \deleted{\unsure{exactement}} précisément, $f$ est une $\mathrm{Hot}_W$-équiv. \unsure{ssi} $\mathrm{hot}_W(f)$ est iso, i.e. $f\in\overline{W}$, le saturé fort de $W$.

Si \uline{$W$ satisfait $W(4b)$}, alors on a également un couple de foncteurs adjoints
\diagram{\begin{tikzcd}[ampersand replacement=\&]
\mathrm{HOT}^{\unsure{\mathrm{lc}|\mathrm{le}}}_W(S')\arrow[r,bend left=8,"f_!^{\mathrm{lc}}"] \& \mathrm{HOT}^{\mathrm{lc}}_W(S)\arrow[l,bend left=11,"f^*_{\mathrm{lc}}"]
\end{tikzcd}}
pour tte flèche $S'\unsure{\mathbin{-}|\to}S$ dans $\mathrm{Cat}$. \deleted{\unsure{On}} On voit aussi dans ce cas que
\margin{$\vert$}
si $f:S'\unsure{\mathbin{-}|\to}S$ est dans $W$, alors $f^*_{\mathrm{lc}}$ et $f_!^{\mathrm{lc}}$ entre $\mathrm{HOT}_W^{\mathrm{lc}}(S)$ et $\mathrm{HOT}_W^{\mathrm{lc}}(S')$ sont des équivalences quasi-inv. l'une de l'autre (moyennant \unsure{$W(4b)$})
% slanted
\margin{\uline{Oubli} : Soit $f:S'\unsure{\mathbin{-}|\to}S$, on a alors $\boxed{f_!\,f^*}:\mathrm{HOT}_W(S)\unsure{\mathbin{-}|\to}\mathrm{HOT}_W(S)$ \deleted{\unsure{est}} : égal à $\mathrm{mult}_{S'/S}$ (cf p.~200) \deleted{\unsure{elle}} Il est aussi égal à $\mathrm{mult}_{\unsure{\xi|\zeta}}$, où $\xi=\mathrm{hot}_{S,W}(S')$, si on suppose ou bien $S'$ $W$-\unsure{lisse|propre} sur $S$, ou bien $S'$ $W$-fibré sur $S$. Cf cas \unsure{gén.} p.~201).}
qui est l'axiome $\mathrm{Der}\,3\unsure{b_2|bz}$ ``version $W$-dérivateur''. On ne peut s'attendre à mieux, du côté $\mathrm{Der}\,3$~!

\textcircled{9} \unsure{Th.} de chgt de base pour $f_!$ dans $\mathrm{HOT}_W$.

Si on a
\diagram{\begin{tikzcd}[ampersand replacement=\&]
X\arrow[d,"{\unsure{t|f}}"'] \& X'\arrow[l,"q"']\arrow[d,"{\unsure{t'|f'}}"]\\
S\& S'\arrow[l,"p"]
\end{tikzcd}}
, avec $f$ $W$-lisse \uline{ou} $\unsure{g|q}$ $W$-propre
